\section{Requirements}
\emph{Show your learning journey through changing requirements.}

Format as stages with clear triggers (The following only serves as an example).
\begin{itemize}
    \item Stage 1: Initial Requirements (Week 1-2). Based on team assumptions and initial research. List core features you initially planned (functional 'what the app does' and non-functional requirements such as usability, accessibility standards). You can use user stories \emph{As a [user type], I want [functionality] so that [benefit]}. 
    \item Stage 2: Post-Interview Requirements (Week 3-4). Based now on user interviews and persona development. Document what changed, e.g. specific requirements that were modified, added, or removed, and also why it changed and how these changes affected your overall solution approach.
    \item Stage 3: Post-Prototype Requirements (Week 5-6). This is now based on usability testing and iterative feedback, where you refine your requirements, resulting in a final list of must-have features.
\end{itemize}